% !TeX root = ../main.tex

\chapter*{Resumen}
\addcontentsline{toc}{chapter}{Resumen}

    Los catálogos automatizados de cúmulos estelares privilegian la homogeneidad 
    sobre la validación dinámica individual, rara vez tratan bajo un mismo 
    procedimiento a cúmulos abiertos y globulares, y el tipo morfológico puede 
    actuar como variable de confusión en los análisis ambientales si no se 
    controla por él. El objetivo general consistió en caracterizar estructural y 
    dinámicamente una muestra de cúmulos abiertos y globulares de la Vía Láctea 
    a partir de Gaia DR3 (catálogo de Hunt y Reffert, 2023), mediante perfiles 
    de King 2D y 3D validados cinemáticamente con la ecuación de Jeans 
    isotrópica, evaluando su relación con la edad, el tiempo de relajación de 
    dos cuerpos y el entorno galáctico. Se ajustó el perfil de densidad de King 
    proyectado (2D) y reconstruido tridimensionalmente (3D) mediante proyección 
    gnomónica y muestreo de rechazo de Monte Carlo, aceptando soluciones con 
    bondad de ajuste y consistencia física entre radios de marea 2D y 3D. La 
    estructura se contrastó con un modelamiento cinemático independiente 
    mediante la ecuación de Jeans y se complementó con el cálculo del tiempo de 
    relajación de dos cuerpos. Del catálogo original (7 167 cúmulos, 1 291 929 
    estrellas miembro) se derivó una muestra final de 100 cúmulos (54 abiertos, 
    46 globulares), seleccionados por riqueza estelar y simetría morfológica 
    suficiente; el ajuste 3D fue aceptado en 59 sistemas.

    La validación frente al catálogo arrojó correlaciones significativas en 
    radio de núcleo ($r = +0.809$), radio de marea ($r = +0.523$) y, 
    en abiertos, masa total ($r = +0.867$); la consistencia cruzada King-Jeans 
    fue igualmente sólida (radio de núcleo: $r = +0.760$; radio de marea: 
    $r = +0.685$), sin discrepancias superiores a un factor de 2 en los 59 
    aceptados. El hallazgo central fue el contraste entre la nula evolución de 
    los parámetros geométricos y la fuerte evolución del estado de relajación 
    dinámica: en la submuestra de 52 cúmulos abiertos con edad disponible, ni el 
    radio de núcleo, ni el radio de marea, ni la concentración correlacionaron 
    con la edad ($|r| \leq 0.17$); la fracción del radio termalizada aumentó con 
    la edad ($r = +0.804$) y el cociente entre el tiempo de relajación al radio 
    de marea y la edad disminuyó con la edad ($r = -0.793$). El 92,3 \% (48 de 
    52) tuvo el núcleo completamente relajado; los 4 restantes tenían, sin 
    excepción, edad menor a 0.4 Ga. El radio galactocéntrico también afecta al 
    radio de núcleo ($r = -0.440$): los cúmulos más internos presentan núcleos 
    más extendidos, consistente con un campo de marea galáctico más intenso. 
    Finalmente, la correlación aparente entre la altura galáctica $|Z|$ y el 
    radio de marea en la muestra completa ($r = +0.531$) resultó un artefacto de 
    mezclar poblaciones morfológicamente distintas: desaparece en abiertos 
    ($r = +0.169$, no significativa) y solo persiste en globulares 
    ($r = +0.664$). En conjunto, estos resultados muestran que combinar el 
    ajuste tridimensional de King con la validación cinemática de Jeans aporta 
    información dinámica no accesible en catálogos de proyección 2D, ofreciendo 
    un marco metodológico replicable para la caracterización de cúmulos 
    estelares con datos de Gaia.